\documentclass{situationist-preprint}

\title{Authenticity as Orientation}
\msid{EXI-08-a}
\msfamily{EXI-08}
\msdiscipline{Meaning, Death and the Absurd}
\mstype{Principal manuscript}
\mskeywords{authenticity; Heidegger; Sartre; bad faith; the \enquote{they}; Taylor; Guignon; orientation}

\begin{document}
\maketitle

\begin{abstract}
Is authenticity possible? The ideal asks that a life be one's own rather than borrowed from the anonymous public that Heidegger calls the \enquote{they}. Sceptics argue that the ideal cannot be met: every self-interpretation is drawn from a shared stock of meanings, and every attempt to coincide with oneself is, by Sartre's analysis, a form of bad faith. We argue that authenticity is possible, as an orientation rather than an achievement. We distinguish two readings of the ideal. On the achievement reading, authenticity is a state in which one's commitments have been fully appropriated and are held as one's own; on the orientation reading, it is a standing direction of concern in which the question whether one's commitments are one's own is kept open and answered in action. We show that the sceptical arguments refute the achievement reading and leave the orientation reading untouched, and that the orientation reading is not trivial, since it has conditions of failure. The result agrees with PHE-07 on authentic being-towards-death and generalizes it.
\end{abstract}

\section{The problem}

The modern ideal of authenticity holds that there is a way of being human that is one's own, and that one ought to live in that way rather than in conformity to what is expected \citep{Trilling1972,Taylor1991}. Its classical philosophical statement is Heidegger's. Everyday existence, on his account, is lived in the mode of the \enquote{they} (\emph{das Man}): one does, judges and enjoys as one does, and the possibilities one takes up are those that the public interpretation of the world makes available \citep[\S 27]{Heidegger1927}. Authenticity (\emph{Eigentlichkeit}) is the mode in which existence takes over its possibilities as its own, and Heidegger ties it to anticipation of death and the call of conscience \citep[\S\S 53--60]{Heidegger1927}. Sartre's version runs through the analysis of bad faith: human reality is not what it is and is what it is not, so that every attempt to be one's role, one's character or one's past is a flight from freedom \citep[Part I, ch.~2]{Sartre1943}.

Both accounts make the ideal look unattainable. Heidegger insists that the \enquote{they} is an existential structure, not a lapse, and that authentic existence is not a detachment from it but a modification of it \citep[\S 27]{Heidegger1927}. Sartre argues that sincerity, the ideal of being what one is, is itself in bad faith, and he defers the description of authenticity, which he calls a self-recovery of corrupted being, to another work \citep[Part I, ch.~2]{Sartre1943}. Critics have drawn the conclusion that the ideal is incoherent or that it is a cultural product posing as a natural one. \citet{Taylor1991} and \citet{Guignon2004} defend the ideal against these critics by arguing that it must be understood dialogically, against horizons of significance that one does not choose. We argue that this defence succeeds only if authenticity is understood as an orientation, and that so understood it is possible.

\section{Two readings}

\begin{definition}[Achievement]\label{d:ach}
An agent is \emph{authentic in the achievement sense} at a time if each of her central commitments has been appropriated by her, in that she has examined it, endorsed it on grounds she recognizes as her own, and holds it in a way that no longer depends on its being what one does.
\end{definition}

\begin{definition}[Orientation]\label{d:or}
An agent is \emph{authentic in the orientation sense} over a period if, throughout it, (i) she treats the question whether her central commitments are her own as open; (ii) she responds to occasions that raise the question, such as conflict, loss or the prospect of death, by examining the commitments concerned rather than by covering the question over; and (iii) her examinations make a difference to what she does.
\end{definition}

Definition~\ref{d:ach} describes a state that, once reached, could be held. Definition~\ref{d:or} describes a direction of concern that has to be sustained: its object is not a condition of the self but a question, and it is satisfied by how the agent stands to the question over time. Clause (iii) excludes idle reflection; clause (ii) is the point at which the orientation can fail.

\begin{principle}[Public sources]\label{p:pub}
Every interpretation an agent can give of her commitments, including the interpretation in which she endorses them as her own, is drawn from meanings that are shared and that she did not originate.
\end{principle}

\begin{principle}[Non-coincidence]\label{p:nonc}
No agent coincides with her commitments: whatever she has endorsed remains, at each later time, something she can take up or reconsider.
\end{principle}

Principle~\ref{p:pub} is Heidegger's thesis that the \enquote{they} is constitutive of existence, and Taylor's thesis that identity is formed dialogically against horizons of significance \citep[ch.~4]{Taylor1991}. Principle~\ref{p:nonc} is the core of Sartre's analysis of bad faith, stripped of its ontology of the for-itself: an agent's past endorsements do not settle her present relation to them.

\section{The argument}

\begin{lemma}\label{l:noach}
Authenticity in the achievement sense is not possible.
\end{lemma}

\begin{proof}
Suppose an agent were authentic in the achievement sense at $t$. Then each central commitment is held in a way that no longer depends on its being what one does. By Principle~\ref{p:pub}, the grounds on which she endorses it are drawn from shared meanings, so her holding it depends on public interpretations, and the appropriation is at most partial. By Principle~\ref{p:nonc}, even a completed appropriation would not be held: at each later time the commitment is again something she can take up or reconsider, so the state Definition~\ref{d:ach} describes would have to be re-entered rather than possessed. A state that cannot be completed and, if completed, cannot be held is not one that an agent can be in.
\end{proof}

\begin{lemma}\label{l:or}
Authenticity in the orientation sense is consistent with Principles~\ref{p:pub} and~\ref{p:nonc}.
\end{lemma}

\begin{proof}
Clause (i) of Definition~\ref{d:or} requires only that the agent treat a question as open, which Principle~\ref{p:nonc} says it is. Clause (ii) requires examination, which proceeds with shared meanings, as Principle~\ref{p:pub} says it must; nothing in the definition asks that the examination escape them. Clause (iii) requires that examination make a difference, which is possible precisely because commitments are not settled once endorsed. The two principles that defeat the achievement reading are among the conditions of the orientation.
\end{proof}

\begin{theorem}\label{t:main}
Authenticity is possible as an orientation, though not as an achievement.
\end{theorem}

\begin{proof}
By Lemma~\ref{l:noach}, it is not possible as an achievement. By Lemma~\ref{l:or}, the orientation is consistent with the facts that rule out the achievement. It remains to show that it is realizable and not trivially so. It is realizable because each clause is a pattern of response that ordinary agents exhibit: a person who, after a bereavement, asks whether the career she pursued was hers or her parents', and acts on the answer, satisfies (ii) and (iii) on that occasion. It is not trivial, because it can fail: an agent who meets every such occasion by retreating to what one does, or who treats the question as answered once and for all, violates (i) or (ii). An orientation with conditions of failure that some agents meet is a possible and substantive form of authenticity.
\end{proof}

The theorem generalizes the result of PHE-07, which argues that authentic being-towards-death is possible as an orientation but not as a stable achievement, because anticipation is constantly drawn back into the public time of the \enquote{they}. Anticipation of death is the occasion Heidegger privileges, but the structure is general: the pull of the \enquote{they} is a standing condition, and authenticity is the standing response to it.

\section{Objections and replies}

\begin{objection}[Deflation]\label{o:defl}
An orientation that never arrives is mere aspiration. To call it authenticity changes the subject.
\end{objection}

\begin{reply}
Many ideals are of this form. Open-mindedness and honesty are not states one reaches and then possesses; they are ways of responding that must be maintained, and an agent who has them can lose them. Heidegger's own text supports the reading: authentic existence is a modification of the \enquote{they}, which remains, and resoluteness must hold itself free for taking back its decision \citep[\S\S 27, 62]{Heidegger1927}. The achievement reading, not the orientation reading, changes the subject.
\end{reply}

\begin{objection}[Social origin]\label{o:soc}
The ideal of authenticity is itself a historical product of modern culture. To pursue it is to do what one does in that culture, which is inauthentic by the ideal's own standard.
\end{objection}

\begin{reply}
The objection relies on the achievement reading, on which authenticity would require commitments with no public source. On the orientation reading, the ideal asks how one holds one's commitments, not where they come from; a commitment of public origin, including the commitment to authenticity, can be examined and owned \citep{Taylor1991,Guignon2004}.
\end{reply}

\begin{objection}[Self-deception]\label{o:sd}
An agent can believe she is keeping the question open while covering it over. Since she cannot be sure which she is doing, the orientation is unknowable and therefore idle.
\end{objection}

\begin{reply}
Clause (iii) makes the orientation partly a matter of conduct, which others and the agent herself can observe over time. The orientation does not require certainty that one has it, and its possibility does not depend on its being known.
\end{reply}

\section{Conclusion}

Authenticity is possible, as an orientation rather than an achievement. The facts the sceptic relies on, that self-interpretation draws on shared meanings and that no agent coincides with her commitments, show that no one can reach a state of fully owned existence and keep it. They do not show that one cannot sustain a direction of concern in which the question of ownership is kept open and answered in action, and this is what authenticity, read as orientation, requires.

\begin{thebibliography}{9}\small
\bibitem[Guignon(2004)]{Guignon2004} Guignon, C. (2004). \emph{On Being Authentic}. Routledge.
\bibitem[Heidegger(1927)]{Heidegger1927} Heidegger, M. (1927). \emph{Sein und Zeit}. Niemeyer. Trans. J. Macquarrie and E. Robinson, \emph{Being and Time}, Harper \& Row, 1962.
\bibitem[Sartre(1943)]{Sartre1943} Sartre, J.-P. (1943). \emph{L'\^etre et le n\'eant}. Gallimard.
\bibitem[Taylor(1991)]{Taylor1991} Taylor, C. (1991). \emph{The Ethics of Authenticity}. Harvard University Press.
\bibitem[Trilling(1972)]{Trilling1972} Trilling, L. (1972). \emph{Sincerity and Authenticity}. Harvard University Press.
\end{thebibliography}

\end{document}
